\section{问题重述及全文大致思路}
\subsection{输入、决策与输出}
对每张计算图，输入包括 Tensor 集合 $T$、核内操作集合 $V$、边集合 $E$ 以及题面固定配置。算法决定三类变量：操作到子图的映射 $g(v)$、子图到核心的映射 $a(s)$ 和每个核心上的子图顺序 $\sigma_k$。最终方案用题面要求的 `node\_to\_subgraph` 和 `core\_schedules` 两个字段输出；操作发射时刻、COPY 插入、换入换出以及 Cache 替换均由评估程序根据场景规则确定。

\begin{table}[H]\centering\caption{逐问交付接口（实验字段待补）}
\begin{tabularx}{\linewidth}{lYYY}\toprule
问题 & 输入与硬约束 & 方法输出 & 待补的正式证据\\\midrule
1 & 计算图、$K=1\ldots5$、场景 A & 子图映射、核心顺序、Makespan 与额外搬运 & 平均加速比曲线、逐用例表\\
2 & 同一计算图、私有 L1/UB、场景 B & 合并 Task 方案、容量合法性、搬运分解 & 加速比、Spill、容量与顺序对照\\
3 & 场景 B 加共享 FIFO L2 & Cache 感知方案、查询与带宽分析 & 无 L2/有 Cache 曲线、命中率、逐用例配对\\\bottomrule
\end{tabularx}\end{table}

\subsection{三问之间的计算链}
三问共享计算图解析、方案合法性检查和核内展开，但 Task 组织不同。问题一的结果不直接作为问题二的数值结论；问题二的成功方案为问题三提供同核 Task 结构和无 L2 的比较入口。后续实验将为每个图、场景和核数保存初始化方案与最终方案的对应关系，避免把不同 Idea 或不同运行的数字拼接在同一表中。

\resultbox{18mm}{输入图、候选方案、官方评估和论文证据的流程图（待绘制方法图）}

\subsection{评价顺序}
对任一场景 $r$，先检查节点覆盖、子图依赖无环和同核顺序，再进行核内拓扑排序、容量处理和多核模拟。算法搜索阶段使用轻量评分减少无效完整调用；正式结论只能来自允许的运行清单和官方评估字段。没有完成的运行只在结果表保留状态和缺失原因，不用零值填补。
